\documentclass[12pt]{article}
\usepackage[utf8]{inputenc}
\usepackage[T1]{fontenc}
\usepackage{hyperref}
\title{Ami Language Resources for K-12 Curriculum Development}
\author{Compiled by Claude Sonnet 5.5 (claude-sonnet-5-5) for Andr\'as Kornai}
\date{}
\begin{document}
\maketitle

\section{Introduction}
Ami (also Emmi; ISO 639-3 amy) is a Western Daly language of the Northern Territory, Australia, traditionally spoken on the coast around Anson Bay. Published estimates are small: about two dozen speakers in the late 1990s (Ford) and about 30 in 2003 (Grimes). Very little material is available online.

\section{Primary Educational Resources}
The Living Archive of Aboriginal Languages holds \emph{Ami Picture Stories} (1987, with parallel English text; by Marjorie and Veronica Bil Bil and Linda Yarrowin): \url{https://territorystories.nt.gov.au/10070/886404/0}

\section{Linguistic and Reference Materials}
\begin{itemize}
\item Lysbeth Ford (1998), ANU PhD thesis, \emph{A description of the Emmi language of the Northern Territory of Australia}, the main grammar. A PDF is reported to be in the ANU repository, but no URL was seen in the results.
\item Grambank entry: \url{https://grambank.clld.org/languages/amii1238}
\item WALS entry: \url{https://wals.info/languoid/lect/wals_code_emm}
\item Wikipedia overview: \url{https://en.wikipedia.org/wiki/Emmiyangal}
\end{itemize}

\section{Religious and Community Resources}
None found.

\section{Digital and Audio Resources}
None found beyond the archive item above.

\section{Curriculum Development Context}
No school programme was found. The grammar and the archive's picture stories are the only identified starting points; any curriculum would depend on collaboration with descendants of Ami speakers in the Wadeye region.
\end{document}